\chapter[Train Track Case Analysis: Track 13]{Train Track Case Analysis: Track 13}
\label{app:case_analysis_track13}

This is Track 13:

\section{Case A: Triangles Fixed}\label{CaseA Track 13}
Let us assume the triangles are fixed. Hence, both $\fbd$ and $\fbdb$ both pass over $d$/$\dbar$ after an even number of letters. We shall call this \textit{Assumption A}. Then $\fbdb$ can begin with $o$, $i$, $s$, or $p$; and $\fbd$ can begin with $r$ or $b$.

\subsection{RNT}\label{RNT}
Suppose that $\fbdb$ began with $o^*$. Then it must be $o^+$:
\[
\fbo = o^+...
\]

Then $\fbd$ can begin with $r^+$ or $b^-$.

Suppose also that $\fbd$ began with $r^*$. Then it must be $r^+$, and then $b^+$:
\[
\fbd = r^+ b^+ d...
\]



This is RNT1.

Here the next letter is either $p^*$ or $s^*$, and there might then be twists with $i$ before going $p^+$/$s^+$ right before crossing $\dbar$. Suppose it were the former case. Then in fact there cannot be such initial twists as this would force $\fbi$ to trace. 

Thus we must have the situation depicted in RNT2.

Here the dotted lines cannot join. Otherwise, $\fbp$ will trace:
\[
\fbp = \dbar b^- r^- d p^\times
\]

This is also true in general if $\fbp$ were to go below $\fbd$. 

On the other hand, if $\fbp$ were to go above $\fbb$, then either $\fbp$ or  $\fbs$ will trace, as shown in RNT3.

Thus, the only possibility is for $\fbp$ and $\fbs$ to go in between $\fbb$ and $\fbd$. 

This is RNT4.

Here $\fbdb$ cannot go in between $\fbb$ and $\fbs$ or it will be in Situation $\beta$. Thus it must be above to give RNT5. But now the dotted lines are in Situation $\alpha$.

This is a contradiction.

Going back to the situation depicted in RNT1. Suppose $\fbdb$ goes down at $s$ before crossing $\dbar$. Then again in factr it cannot have any twists with $i$ or $\fbi$ will trace.

This gives us RNT6.

Here the situation is similar to the previous scenario. If $\fbp$ were to go below $\fbd$ then $\fbs$ will trace, as shown in RNT7. If $\fbp$ were to go above $\fbb$ then again $\fbs$ traces, as shown in RNT8. Thus we must have the situation depicted in RNT9, where the dotted lines again are in Situation $\alpha$.

This is a contradiction.

Back to the beginning of RNT \ref{RNT}. Suppose $\fbd$ were to begin with $b^-$:
\[
\fbd = b^- r^- d...
\]

This is RNT10.

By the same reasoning as before, $\fbdb$ can either traverse:
\[
\fbdb = o^+ p^+ \dbar...
\]
or
\[
\fbdb = o^+ s^+ \dbar...
\]

\subsubsection{Bort}
Suppose it were the former case:
\[
\fbdb = o^+ p^+ \dbar...
\]

This is Bort1.

Here, by the same reasoning as before, $\fbp$ and $\fbs$ can only go in between $\fbd$ and $\fbr$, giving us Bort2.

Consider $\fbdb$. If it went below $\fbd$, then it would be in Situation $\beta$. If it went above $\fbd$, then it would violate Rule F1 after crossing $d$ to the right-hand side due to the position of $\fbr$. Thus the dotted lines must join. On the other hand, $\fbs$ must continue with $r^+$ and to the left of $\fbb$ or else $\fbb$ will trace immediately after crossing $d$; and $\fbp$ must continue with $r^+$ and to the left of $\fbb$ as well, or else $\fbb$ will be absorbed between $\fbs$ and $\fbp$.

Thus we have Bort3.

Here $\fbs$ must continue with $b^+$ and to the left of $\fbb$. Otherwise, $\fbp$ would be forced to continue with $b^-$ and to the right of $\fbs$, and then trace
\[
\fbp = r^+ \hat{b}^- r^- d p^\times
\]

Thus we have
\[
\fbs = \dbar r^+ \hat{b}^+ r^*...
\]

which forces $\fbb$ to trace:
\[
\fbb = r^+ b^\times
\]

This is a contradiction.


\subsubsection{Cess}
Back to the end of RNT \ref{RNT}. Suppose 
\[
\fbdb = o^+ s^+ \dbar...
\]

This is Cess1.

Again for the same reasoning as before, $\fbs$ and $\fbp$ must go in between $\fbd$ and $\fbr$. But now $\fbr$ will trace:
\[
\fbr = \dbar o^+ s^+ \dbar r^\times
\]
as shown in Cess2.

\subsection{PWK}\label{PWK}
Back to the end of Case A \ref{CaseA Track 13}. Suppose $\fbdb$ began with $i$, in which case it must be $i^+$ as to not violate Assumption A; and suppose $\fbd$ began with $r$.

Thus we have the situation depicted in PWK1.

Here $\fbdb$ can either continue with $p^+ \dbar$ or $s^+ \dbar$.

\subsubsection{Dowl}
Suppose
\[
\fbdb = i^+ p^+ \dbar...
\]

This is Dowl1.

Here $\fbb$ cannot go above $\fbdb$ as it would trace:
\[
\fbb = \dbar p^- i^- \dbar b^\times
\]

So either $\fbb$ goes in between $\fbdb$ and $\fbs$, or it goes below $\fbp$. If it were the latter, then $\fbd$ continues:
\[
\fbd = r^+ b^+ d i^+...
\]
here the next letter cannot be $s^*$, as then it would violate Rule F1. So we have
\[
\fbd = r^+ b^+ d i^+ p^+ \dbar...
\]

This is Dowl2.

Here the dotted lines are in Situation $\alpha$. 

So suppose it were the former case, that $\fbb$ went below $\fbp$. But then again we have
\[
\fbd = r^+ b^+ d i^+ p^+ \dbar...
\]
and now this forces $\fbb$ to trace, as shown in Dowl3.

This is a contradiction.

\subsubsection{Fubs}
Back to the end of PWK \ref{PWK}. Suppose $\fbdb$ continued with $s^+ \dbar$:
\[
\fbdb = i^+ s^+ \dbar...
\]

This is Fubs1.

Here if $\fbd$ went below $\fbp$, then it would cause the dotted lines to be in Situation $\alpha$, as shown in Fubs2.

If instead the dotted lines join here, then the following is forced for $\fbb$:
\[
\fbb = d p^*...
\]
as the alternative would cause it to trace:
\[
\fbb = d s^- i^- \dbar b^\times
\]

But now $\fbp$ will trace:
\[
\fbp = \dbar b^- r^- p^\times
\]
as shown in Fubs2.

This is the same for the case where $\fbd$ went above $\fbdb$, as shown in Fubs3.

This is a contradiction.

\subsection{BAF}\label{BAF}
Back to the end of Case A \ref{CaseA Track 13}. Suppose $\fbdb$ began with $i$, in which case it must be $i^+$ as to not violate Assumption A; and suppose $\fbd$ began with $b^-$.

Thus we have the situation depicted in BAF1.

Here again $\fbdb$ can contiue with
\[
\fbdb = i^+ p^+ \dbar...
\]
or
\[
\fbdb = i^+ s^+ \dbar...
\]

\subsubsection{Grig}\label{Grig}
Suppose
\[
\fbdb = i^+ p^+ \dbar...
\]

This is Grig1.

Here $\fbs$ and $\fbp$ cannot go under $\fbr$. Otherwise, $\fbp$ will trace:
\[
\fbp  \dbar b^- r^- d p^\times
\]

On the hand, if $\fbp$ and $\fbs$ went above $\fbd$, then we have
\[
\fbd = b^- r^- d i^+ p^+...
\]
and now the dotted lines are in Situation $\alpha$, as shown in Grig2.

Thus $\fbs$ and $\fbp$ must go in between $\fbs$ and $\fbp$:
\[
\fbp = \dbar r^*...
\]
and
\[
\fbs \dbar r^*...
\]

This is Grig3.

Here $\fbr$ must end at $o^\circ$. Otherwise it would trace:
\[
\fbr = d \hat{o}^+ \dbar r^\times
\]
or
\[
\fbr = d o^- \dbar b^- r^\times
\]

Here the dotted lines must join. Otherwise, if $\fbd$ went below $\fbdb$ then $\fbd$ immediately enters Situation $\beta$. If $\fbd$ went above $\fbdb$, then $\fbdb$ immediately enters Situation $\beta$.

This gives us the situation depicted in Grig4.

Here $\fbb$ can begin and end with $r^\circ$, or begin and continue with $b^-$.
\subsubsubsection{Chirk}\label{Chirk}
Suppose
\[
\fbb = r^\circ
\]
Then we have Chirk1.

Here $\fbs$ and $\fbo$ must continue with $b^+$ and to the left of $\fbp$. Otherwise, $\fbp$ would be forced to continue with $b^-$, as it either traces:
\[
\fbp = \dbar r^+ \hat{b}^- d p^\times
\]
or is absorbed:
\[
\fbp = \dbar r^+ \hat{b}^- d \Abso{\fbs}{\fbp}
\]

Then $\fbo$ continues:
\[
\fbo = p^+ \dbar r^+ \hat{b}^+ r^- d i^+ s^*...
\]
and $\fbs$ continues:
\[
\fbs = p^+ \dbar r^+ \hat{b}^+ r^- d i^+ p^+ \dbar r^+ b^+ d o^+ \dbar r^+ b^- r^- d p^- i^*...
\]

This is Chirk2.

Here $\fbi$ must begin and end at $p^\circ$. Otherwise, it would trace:
\[
\fbi = p^+ \dbar r^+ b^+ r^- d i^\times
\]

Now consider $\fbp$. It must end at $b^\circ$. Otherwise, it would be absorbed:
\[
\fbp = ... b^- r^- d p^- \Abso{\fbi}{ \fbo}
\]
or trace:
\[
\fbp = ... b^+ r^- d i^+ p^\times
\]

Now consider $\fbo$ and $\fbs$. They both must end. Indeed, otherwise $\fbs$ will trace:
\[
\fbs=...i^- r^+ \dbar r^+ b^+ r^- d i^+ s^\times
\]
or be absorbed:
\[
\fbs=...i^- r^+ \dbar r^+ b^+ r^- d o^- \dbar r^- d p^- i^- \dbar r^+ b^-r^- d p^- \Abso{\fbi}{\fbo}
\]
or
\[
\fbs = ... i^+ p^+\dbar r^+ b^+ r^- d o^- \dbar b^- r^- d p^- i^- \dbar r^+ b^- r^- d \Abso{\fbs}{\fbp}
\]

As soon as $\fbs$ ends, if $\fbo$ did not also end, it would be absorbed:
\[
\fbo =...s^+ i^- \dbar r^+ b^- r^- d p^- \Abso{\fbi}{\fbo}
\]
or trace:
\[
\fbo =...s^- \dbar r^+ b^- r^- d p^- i^+ p^+ \dbar r^+ b^+ r^- d o^\times
\]

Thus we have Chirk3, which is Candidate 1.

\subsubsubsection{Dwalm}\label{Dwalm}
Back to the end of Grig \ref{Grig}. Suppose
Suppose
\[
\fbb = r^- d...
\]
This is Dwalm1.

Here $\fbb$ can go into 1, 2, or 3.

\subsubsubsubsection{Bocage}
Suppose $\fbb$ went into 1. Then it continues:
\[
\fbb = r^- d i^+ s^*...
\]

This is Bocage1.

Here $\fbs$ and $\fbo$ both must continue with $b^=$ and to the left of $\fbp$. Otherwise, $\fbp$ would be forced to continue with $b^-$ and be absorbed:
\[
\fbp = \dbar r^+ \hat{b}^- d \Abso{\fbs}{\fbp}
\]
or trace:
\[
\fbp = \dbar r^+ \hat{b}- d p^\times
\]

Then $\fbo$ continues:
\[
\fbo = p^+ \dbar r^+ \hat{b}^+ r^- d i^+ s^*...
\]

and $\fbs$ continues:
\[
\fbs = \dbar r^+ \hat{b}^+ r^- d i^+ p^+ \dbar r^+ b^+ d o^+ \dbar r^*...
\]

This is Bocage2. Here it is clear that $\fbp$ and $\fbi$ must end immediately, otherwise they would trace.

Consider $\fbs$. If it continued with $r^+$, it would be absorbed:
\[
\fbs = ...r^+ d o^- \dbar b^- r^- d p^- i^- \dbar r^+ b^- r^- d \Abso{\fbs}{\fbp}
\]
While it it continued with $r^-$ it would trace:
\[
\fbs = ...r^- d i^+ s^\times
\]

Thus it must end here.

This is Bocage3.

Here if $\fbo$ continued with $s^\circ$ or $s^+$, forcing $\fbb$ to continue with $s^+$, then $\fbb$ would trace:
\[
\fbb = ...s^+ i^- \dbar r^+ b^\times
\]

This means $\fbo$ must continue with $s^-$ and to the right of $\fbb$. But then it would also trace:
\[
\fbo = ...s^- \dbar r^+ d o^\times
\]

This is a contradiction.

\subsubsubsubsection{Cagmag}\label{Cagmag}
Back to the end of \ref{Dwalm}. Suppose $\fbb$ went into 2:
\[
\fbb = r^- d p^*...
\]

Then this immediately forces:
\[
\fbp = p^+ \dbar r^+ b^*...
\]

This is Cagmag1.

Here $\fbo$ must continue with $b^+$ and to the left of $\fbi$. Otherwise, $\fbi$ would continue with $b^-$ and trace:
\[
\fbi = ...b^- r^- d p^- i^\times
\]

Thus we have
\[
\fbo = p^+ d r^+ b^+ r^- d p^*...
\]

This is Cagmag2.

Here $\fbs$ cannot continue with $r^\circ$ or $r^-$, as that would force $\fbp$ to continue with $r^-$ and then immediately trace after crossing $d$. So suppose $\fbs$ continued with $r^+$ and to the left of $\fbp$. Then it continues:
\[
\fbs = \dbar r^+ d...
\]

This is Cagmag3. 

Here $\fbs$ can either continue with $o^-$ or $i^+$. Suppose it were the former:
\[
\fbs = ...o^- \dbar b^- r^- d p^- i^- \dbar r^*...
\]

This is Cagmag4.

Here $\fbo$ cannot continue with $p^\circ$ or $p^+$ as that would force $\fbb$ to continue with $p^+$, and thus trace:
\[
\fbb = r^- d \hat{p}^+ \dbar r^+ b^*...
\]

However, if $\fbo$ continued with $p^-$ and to the right of $\fbb$, then it also traces:
\[
\fbo = ...p^- \dbar r^+ d o^\times
\]

So suoopose it were the latter case, that $\fbs$ continued with $i^+$:
\[
\fbs = ...i^+ p^+ \dbar r^+ b^+ d o^+ \dbar r^*...
\]

Here $\fbp$ it is clear must end at $r^\circ$ or it will trace. So this forces the next lette for $\fbs$ to be $r^-$:
\[
\fbs = ...i^+ p^+ \dbar r^+ b^+ d o^+ \dbar \hat{r}^- d p^*...
\]

This is the situation depicted in Cagmag5.
Here $\fbo$ must continue with $p^-$. Otherwise, it would force $\fbb$ to continue with $p^+$ and immediately trace after crossing $\dbar$. Now if $\fbo$ went in between $\fbb$ and $\fbs$ then it would trace:
\[
\fbo = ...p^- \dbar r^+ d o^\times
\]

So it must go in between $\fbs$ and $\fbi$. 

This is Cagmag6. 

Here $\fbi$ must end. Indeed, otherwise it trace:
\[
\fbi = ...b^+ r^- d p^- \dbar r^+ d i^\times
\]
or
\[
\fbi = ...b^+ r^- d p^- \dbar r^+ d o^- \dbar b^- r^- d p^- i^\times
\]
or it it absorbed:
\[
\fbi = ...b^- r^- d p^- \Abso{\fbi}{\fbs}
\]


Thus we have Cagmag7.

Here we claim that $\fbb$ must end at $p^\circ$. Indeed, otherwise we either have
\[
\fbb = r^- d \hat{p}^+ \dbar r^+ b^\times
\]
or it continues with $p^+$, which is the sitation depicted in Cagmag8. But here $\fbo$, $\fbb$, and $\fbs$ are stuck in a three-way clockwise path mill between $s$ and $p$. This is a contradiction.

Thus we must have
\[
\fbb = r^- d p^\circ
\]

This is Cagmag9.

Here we claim that both $\fbs$ and $\fbo$ must end here. Indeed, consider $\fbs$. If it did not end here, then
\[
\fbs = ...i^+ p^+ \dbar r^+ b^+ r^- d p^- \dbar r^+ d o^- \dbar b^- r^- d p^- i^- \dbar r^- d \Abso{\fbs}{\fbp}
\]
or
\[
\fbs = ...i^- p^+ \dbar r^+ b^+ r^- d p^+ \dbar r^+ b^- r^- d p^- \Abso{\fbi}{\fbo}
\]

Thus $\fbs$ must end at $i^\circ$. This then also force $\fbo$ to end at $s^\circ$, otherwise it would eventually trace.

This is Cagmag10, which is Candidate 2.

\subsubsubsubsection{Boofus}\label{Boofus}
Back to the end of \ref{Dwalm}. Suppose $\fbb$ went into 3:
\[
\fbb = r^- d p^*...
\]

This is Boofus1.

Here $\fbo$ must continue with $r^+$ and to the left of $\fbp$. Otherwise, $\fbp$ would immediately trace after crossing $d$. The same is true for $\fbs$. The two continue to give the situation depicted in Boofus2.

Here we claim that $\fbp$ must end at $r^\circ$. Indeed, otherwise it would trace:
\[
\fbp = \dbar \hat{r}^+ d p^\times
\]
or absorbed:
\[
\fbp = \dbar \hat{r}^- d \Abso{\fbs}{\fbp}
\]

Thus 
\[
\fbp = \dbar r^\circ
\]

now consider $\fbi$. It must end immediately at$p^\circ$. Otherwise, it would trace:
\[
\fbi = p^+ \dbar r^+ d i^\times
\]

This is Boofus3.

Here $\fbs$ must continue with $i^+$ and to the left of $\fbb$. Otherwise, $\fbb$ would continue with $i^-$ and to the right of $\fbs$. In which case, we could have the situation depicted in Boofus4.

Here $\fbs$ must end or it will eventually trace. Whereas if $\fbo$ continued with $s^\circ$ or $s^+$, then $\fbb$ would eventually be absorbed between $\fbi$ and $\fbo$; but if $\fbo$ continued with $s^-$, then it would eventually trace. 

Thus we must have the situation depicted in Boofus5, where $\fbs$ continued with $i^+$ and $\fbb$ must end or it will eventually trace. Here $\fbo$ must end as well or it will eventually trace; and the same is true for $\fbs$. This gives Candidate 3, shown in Boofus6.

\subsubsection{Harl}
Back to the end of BAF \ref{BAF}. Suppose 
\[
\fbdb = i^+ s^+ \dbar...
\]

This is Harl1.

For the same reasoning as that in Grig \ref{Grig}, $\fbs$ and $\fbp$ cannot go under $\fbr$. 

On the other hand, if $\fbp$ and $\fbs$ went above $\fbd$, then we would similarly have the situation depicted in Harl2, where  the dotted lines are in Situation $\alpha$.

So the only possibility is for $\fbs$ and $\fbp$ to go in between $\fbd$ and $\fbr$.

This gives Harl3.

If if the dotted lines do not join, then at least one of the two ($\fbd$ or $\fbdb$ will immediately be in Situation $\beta$). Thus the dotted lines must join.

Now $\fbo$ continues:
\[
\fbo = s^+ \dbar r^+ b^+...
\]

This is Harl4. 

Here if $\fbb$ were to end at $r^\circ$, then $\fbo$ is either absorbed:
\[
\fbo = s^+ \dbar r^+ b^+ \hat{r}^- d \Abso{\fbs}{\fbp}
\]
or traces:
\[
\fbo = s^+ \dbar r^+ b^+ \hat{r}^- d o^\times
\]
or
\[
\fbo = s^+ \dbar r^+ b^+ \hat{r}^- d i^+ s^+ \dbar r^+ b^+ d o^\times
\]

Thus we must have
\[
\fbb = r^- d...
\]

and in between itself and $\fbp$, otherwise it would be immediately absorbed between $\fbp$ and $\fbs$, as shown in Harl5. Here notice that regardless of whether the next letter is $o^-$ or $i^+$, it would eventually trace. This is a contradiction.

\subsection{SEN}\label{SEN}
Back to the end of Case A \ref{CaseA Track 13}. Suppose $\fbdb$ began with $s$, in which case it must be $s^-$ as to not violate Assumption A; and suppose $\fbd$ began with $r^+$.

Thus we have the situation depicted in SEN1.

Here we claim the dotted lines must join. Indeed, otherwise, either $\fbd$ or $\fbdb$ immediately enter Situation $\beta$, or the dotted lines enter Situation $\alpha$.

Now consider $\fbp$. It must begin with $i^-$ so as to not cause $\fbs$ to immediately trace. Then it continues:
\[
\fbp = i^- \dbar b^*...
\]
where it must do so going below $\fbb$ in the middle, otherwise this forces $\fbo$ and $\fbi$ to continue with $...r^+b^+ d...$ and thereby causing at least one of them to trace. For the same reason, $\fbb$ must go above $\fbo$ and $\fbi$. This gives us the situation depicted in SEN2.

Here $\fbi$ can continue with $b^\circ$ or $b^-$.

\subsubsection{Marl}
Suppose it were the former case:
\[
\fbi = \dbar \hat{b}^\circ
\]

This is Marl1.

Here $\fbp$ cannot continue in between $\fbi$ and $\fbo$, or it would be absorbed:
\[
\fbp = i^- \dbar \hat{b} d \Abso{\fbo}{\fbi}
\]

So suppose it went in between $\fbr$ and $\fbo$. Then it continues:
\[
\fbp = i^- \dbar \hat{b} d s^- i^- \dbar b^- r^- p^\times
\]
as shown in Marl2.

So suppose it went in between $\fbd$ and $\fbr$. Then $\fbr$ and $\fbo$ continue with $...di^*...$, forcing
\[
\fbs = i^- \dbar b^- r^*...
\]

as shown in Marl3.

Here $\fbp$ must continue with $r^+$ and to the left of $\fbs$. Otherwise, $\fbs$ would trace after crossing $d$. This gives Marl4.

Here if $\fbr$ continued with $i^\circ$ or $i^+$ then $\fbp$ would either be absorbed between $\fbo$ and $\fbi$, or it would escape to the outer perimeter and eventually trace. But if $\fbr$ continued instead with $r^-$, it would trace itself:
\[
\fbr = ...i^- \dbar b^- r^\times
\]

This is a contradiction.

\subsubsection{Peen}\label{Peen}
Back to the end of SEN \ref{SEN}. Suppose $\fbi$ continued with $b^-$.

This is Peen1.

Here $\fbi$ can either go into 1 or 2. Note that it cannot go in between $\fbi$ and $\fbo$ or it would be immediately absorbed. Note also that $\fbp$ can only continue with $r^*$ otherwise it would escapte to the outer perimeter and eventually trace.

\subsubsubsection{Howff}
Suppose $\fbi$ went into 1. Then it continues:
\[
\fbi = \dbar b^- d p^*...
\]

and $\fbr$ continues:
\[
\fbr = b^+ d i^*...
\]

which forces
\[
\fbs = i^- \dbar b^- r^*...
\]
and
\[
\fbp = i^- \dbar b^- r^*...
\]

This is Howff1.

Here $\fbi$ must continue with $p^+$ and to the left of $\fbb$. Otherwise, $\fbb$ would immediately trace after crossing $\dbar$. This is Howff2.

Here both $\fbb$ and $\fbr$ must end. Otherwise, both would eventually trace.

This is Howff3.

Now if $\fbp$ continued with $r^\circ$ or $r^+$, then $\fbs$ would continue with $r^+$ and trace:
\[
\fbs = ...r^+ b^+ d i^+ s^\times
\]

But then if $\fbp$ continued with $r^-$, it would trace:
\[
\fbp = ...r^- b^+ d i^+ \dbar b^- d p^\times
\]

This is a contradiction.

\subsubsubsection{Jasey}\label{Jasey}
Back to the end of Peen \ref{Peen}. Suppose $\fbi$ went into 2. 

This gives Jasey1.

Here $\fbr$ can either end immediately at $b^\circ$, or continue with $b^+$. 

\subsubsubsubsection{Gimmer}
Suppose 
\[
\fbr = b^\circ
\]

This gives Gimmer1.

Here $\fbi$ can either end at $r^\circ$ or continue with $r^-$. Suppose it were the former. Then $\fbp$ will continue with $r^-$ and in between $\fbo$ and $\fbd$ (going in tween $\fbi$ and $\fbo$ would cause it to be absorbed). Then it continues and traces:
\[
\fbp = ...r^- b^+ d s^- i^- \dbar b^- r^- d p^\times
\]
as shown in Gimmer2.

So suppose it were the latter, that $\fbi$ continued with $r^-$, as shown in Gimmer3. But here the behavior of $\fbp$ is identical to the former case. This is a contradiction.

\subsubsubsubsection{Hantle}
Back to the end of Jasey \ref{Jasey}. Suppose $\fbr$ continued with $b^+$. 

This is shown in Hantle1.

Here notice that $\fbr$ will escape to the outer perimeter and eventually trace.

This is a contradiction.

\subsection{HUN}\label{HUN}
Back to the end of Case A \ref{CaseA Track 13}. Suppose $\fbdb$ began with $s$, in which case it must be $s^-$ as to not violate Assumption A; and suppose $\fbd$ began with $b^-$.

Thus we have the situation depicted in HUN1.

Here if the dotted lines were to join, then $\fbi$ or $\fbo$ will eventually trace. For the same reason, the only alternative is for $\fbd$ to go above $\fbo$ and $\fbi$. But now we will have Situation $\alpha$, as depicted in HUN2.

This is a contradiction.

\section{Case B,C,D: Triangles Swapped}\label{CaseBCD Track 13}

Now suppose the triangles are swapped. Then at least one of $\fbd$ and $\fbdb$ must go over an odd number of letters before crossing. We first assume this to be the case for $\fbd$, that is, $\fbd$ crosses an odd number of letters before crossing $\dbar$.

Then $\fbd$ can begin with $o$, $i$, $s$, or $p$; and $\fbdb$ can begin with $r$, $b$, or $d$.

\subsection{GUW}\label{GUW}
Suppose $\fbd$ began with $o$, and $\fbdb$ began with $r$. We remind the reader that we are assuming $\fbd$ must go over an odd number of letters before crossing. This assumption necessitates that eventually $\fbd$ must traverse $o^-$ right before crossing $\dbar$.

Now $\fbb$ must begin with
\[
\fbb = \dbar r^*...
\]
which forces
\[
\fbp = d o^*...
\]
and
\[
\fbs = d o^*...
\]

This is GUW1.

Here notice that $\fbd$ cannot have any initial twists before traversing $o^- \dbar...$, as that would cause either $\fbp$ or $\fbs$ to trace.

Thus we must have GUW2.

Here both dotted lines must continue with $r^-$ otherwise $\fbb$ would immediately trace. But now the dotted lines are in Situation $\alpha$, as shown in GUW3.

This is a contradiction.

\subsection{CPX}\label{CPX}
Back to the end of Case B,C,D \ref{CaseBCD Track 13}.
Suppose $\fbd$ began with $o$, and $\fbdb$ began with $r$. We remind the reader that we are assuming $\fbd$ must go over an odd number of letters before crossing. This assumption necessitates that eventually $\fbd$ must traverse $o^-$ right before crossing $\dbar$.

Again we must have
\[
\fbb = \dbar r^*...
\]
which forces $\fbo$ to trace:
\[
\fbo = d o^\times
\]
as shown in CPX1.

\subsection{SCC}\label{SCC}
Back to the end of Case B,C,D \ref{CaseBCD Track 13}.
Suppose $\fbd$ began with $o$, and $\fbdb$ began with $d$. We remind the reader that we are assuming $\fbd$ must go over an odd number of letters before crossing. This assumption necessitates that eventually $\fbd$ must traverse $o^-$ right before crossing $\dbar$.

We have the situation depicted in SCC1.

Now $\fbo$ and $\fbs$ must continue with $r^-$ and to the right of $\fbb$. Otherwise $\fbb$ would immediately trace.

Here the dotted lines must join. Otherwise, either $\fbd$ went below $\fdbd$, which would put $\fbdb$ in Situation $\beta$; or $\fbd$ went above $\fbdb$, which would put $\fbd$ in Situation $\beta$. Thus they must join.

This is SCC2.

Here $\fbr$ can begin with $p^\circ$, $s^\circ$, $p^-$, or $s^-$.

Suppose $\fbr$ began with $p^\circ$. Then $\fbs$ would trace, as shown in SCC3.

Suppose $\fbr$ began with $s^\circ$. This is SCC4.

Here $\fbb$ must continue with $r^-$ or $\fbs$ would trace immediately after crossing $d$. 

This is SCC5.

Here $\fbi$ can either end at $p^\circ$ or continue with $p^+$ or $p^-$.

Suppose it ended here with $p^\circ$:
\[
\fbi = r^- d p^\circ
\]

Then this is SCC6. Here $\fbo$ must continue with $i^-$ and to the right of $\fbb$. Otherwise, $\fbb$ would be absrobed:
\[
\fbb = ... i^+ s^+ \dbar r^+ \Abso{\fbi}{\fbo}
\]

This gives us SCC7. 

Here if $\fbs$ were to continue with $r^\circ$ or $r^+$, then $\fbo$ would continue with $r^+$ and then trace:
\[
\fbo = ...r^+ d s^- o^\times
\]

But if $\fbs$ continued with $r^-$ and to the right of $\fbo$, then it would also trace:
\[
\fbs = ...r^- d s^\times
\]

This is a contradiction.

So suppose in the situation depicted in SCC7, $\fbi$ continued with $p^+$. Then it continues:
\[
\fbi = ...p^+ \dbar r^+ b^+ d o^+ s^+ \dbar...
\]

This is SCC8. Here if $\fbi$ went above $\fbo$ then $\fbo$ would trace. The same is true if $\fbi$ went be tween $\fbo$ and $\fbb$. Thus $\fbi$ can only go below $\fbb$. This gives SCC9.

Here $\fbo$ must continue with $i^-$ and to the right of $\fbb$. Otherwise, $\fbb$ would be absorbed:
\[
\fbb = ...i^+ s^+ \dbar r^+ \Abso{\fbi}{\fbo}
\]

Thus $\fbo$ continues:
\[
\fbo = ...i^- d^+ \dbar r^*...
\]

as shown in SCC10.

Here $\fbo$ cannot continue with $r^+$, or it will trace after crossing $d$. Thus $\fbs$ and $\fbi$ need to continue with $r^-$ and to the right of $\fbo$. But then $\fbs$ traces:
\[
\fbs = ...r^- d s^\times
\]

So suppose in the situation depicted in SCC7, $\fbi$ continued with $p^-$, as in SCC11. Again here it must go below $\fbo$ to not be absorbed eventually between itself and $\fbo$

Then it continues:
\[
\fbi = ...p^- \dbar r^+ d o^*...
\]

This is SCC12.

Here notice that $\fbs$ must end immediately at $r^\circ$. Otherwise, it would immediately trace after crossing $d$. Now we claim that all of $\fbb$, $\fbi$, $\fbo$, and $\fbp$ must end immediately.

Suppose by contradiction that $\fbb$ did not end. Then suppose it continued with $i^-$, then it will trace:
\[
\fbb =...i^- p^+ \dbar r^+ d s^- o^- \dbar b^\times
\]
as shown in SCC13.

So suppose it continued with $i^+$, then it continues:
\[
\fbb = ...i^+ s^+ \dbar r^+ d o^*...
\]

This is SCC14.

Here if $\fbb$ continued with $o^\circ$ or $o^+$, then $\fbi$ continues with $o^+$ and to the left of $\fbb$. But then it will trace:
\[
\fbi = ...o^+ \dbar r^- d s^- i^\times
\]

But then if $\fbb$ continued with $o^-$ and to the right of $\fbo$, then it continues to give the situation depicted in SCC15.

Here if $\fbb$ continued with $p^+$ then it will immediately be absorbed right after crossing $\dbar$. But then this means $\fbo$ must continue with $p^-$ and to the right of $\fbb$, in which case it would trace:
\[
\fbo = p^- \dbar r^+ d o^\times
\]

This establishes that $\fbb$ must end at $i^\circ$. This is SCC16.

Here let us suppose by contradiction that $\fbi$ did not end. Then suppose it continued with $o^+$:
\[
\fbi = ...o^+ \dbar r^- d p^+ \dbar r^+ b^+ d o^+ s^+ \dbar r^+ d s^- o^- \dbar b^*...
\]

This is SCC17.

But here $\fbp$ is forced to begin with $b^+$, and trace:
\[
\fbp = b^+ d o^+ s^+ \dbar r^- d s^- o^- \dbar b^- r^- d p^\times
\]

This is a contradiction. So suppose instead $\fbi$ continued with $o^-$:
\[
\fbi =...o^- \dbar r^- d p^*...
\]

This is SCC18.

Here if $\fbo$ continued with $p^\circ$ or $p^+$, then $fbi$ would continue with $p^+$ and to the left of $\fbo$, but then it would immediately be absorbed after crossing $\fbdb$. But hen if $\fbo$ continued with $p^-$ and to the right of $\fbi$, then it would trace immediately after crossing $d$.

This establishes that $\fbi$ must end at $o^\circ$.

This is SCC19.

Now suppose by contradiction that $\fbo$ did not end. Then it could only continue with $p^-$:
\[
\fbo = ...p^- \dbar r^+ d o^+ \dbar r^- d s^- i^- s^+ \dbar r^+ d s^- o^- \dbar b^*...
\]

This is SCC20.

Here this forces $\fbp$ to begin with $b^+$, and thereby eventually trace, similar to a previous situation.

This establishes that $\fbo$ must end at $p^\circ$. 

This is SCC21.

Here if by contradiction $fbp$ did not end at $b^\circ$, then
\[
\fbp = b^+ d o^+ s^+ \dbar r^- d s^- i^- s^+ \dbar r^+ d o^- \dbar r^- d p^\times
\]

Thus we must have
\[
\fbp = b^\circ
\]

This establishes Candidate 4, shown in SCC22.

Now back to the Situation depicted in SCC2. Suppose $\fbr$ began with $p^-$:
\[
\fbr = p^-...
\]

But this forces $\fbi$ to trace, as shown in SCC23.

So suppose $\fbr$ began with $s^-$:
\[
\fbr = s^-...
\]
Then $\fbd$ cannot have any initial twists. On the other hand, $\fbi$ must continue with $p^*...$:
\[
\fbi = r^- d p^*...
\]

Here $\fbs$ cannot continue with $b^+$ or it would trace after crossing:
\[
\fbs = b^+ d o^+ s^\times
\]

This means that $\fbp$ must begin with $b^-$ and to the right of $\fbs$, but then it will trace:
\[
\dbp = b^- r^- d p^\times
\]

This is shown in SCC24.

\subsection{LKE}\label{LKE}
Back to the end of Case B,C,D \ref{CaseBCD Track 13}.
Suppose $\fbd$ began with $i$, and $\fbdb$ began with $r$. We remind the reader that we are assuming $\fbd$ must go over an odd number of letters before crossing. This assumption necessitates that eventually $\fbd$ must traverse $i^-$ right before crossing $\dbar$.

We have the situation depicted in LKE1.

Here the dotted lines must join. Otherwise, either $\fbd$ or $\fbdb$ would immediately enter Situation $\beta$.

This is LKE2.

Here $\fbr$ cannot begin with $p^\circ$ or $p^-$, as in that case $\fbo$ or $\fbs$ would immediately trace after crossing and after $p^-$. On the other hand, $\fbr$ also cannot begin with $p^+$ or $s^+$ as then it would immediately trace after crossing $\dbar$.

So it can only begin with $s^\circ$ or $s^-$. 

Suppose it were the former. Then this forces $\fbd$ to not have any initial twists. Also, $\fbi$ must continue with $p^*$, otherwise it would trace. This then forces $\fbs$ to traverse:
\[
\fbs = b^+ d o^*...
\]
otherwise $\fbp$ would begin with $b^-$, and immediately trace after crossing.

This is LKE3.

Here $\fbo$ can either continue with $p^*$, or with $s^-$. Suppose it were the former. Then it cannot continue with $p^+$, or it would trace:
\[
\fbo =...p^+ \dbar r^+ b^+ d o^\times
\]
This means $\fbi$ must continue with $p^-$ and to the right of $\fbo$, and continue:
\[
\fbi = r^- p^- \dbar r^*...
\]

This is LKE4.

Here $\fbb$, $\fbo$ , and $\fbi$ are stuck in a three-way clockwise path mill between $r$ and $p$.

So suppose $\fbo$ continued with $s^-$. Then it continues:
\[
\fbo = r^- d \hat{s}^- i^- \dbar b^*...
\]

forcing
\[
\fbp = b^+ d...
\]

This is LKE5.

Here if the next letter for $\fbp$ were $o^*$, then $\fbs$ must continue with $o^+$. Otherwise, $\fbp$ would continue with $o^-$ and eventually trace. So we have
\[
\fbs = b^+ d \hat{o}^+ \dbar b^*...
\]

This is LKE6. Here $\fbs$, $\fbo$, and $\fbp$ are stuck in a three-way counterclockwise path mill between $b$ and $o$.

So suppose the next letter for $\fbp$ were $i$, then it would be absorbed:
\[ 
\fbp = b^+ d \hat{i}^+ s^+ \dbar r^+ \Abso{\fbi}{\fbo}
\]
as shown in LKE7.

Going back to the situation depicted in LKE2. Suppose $\fbr$ began with $s^-$. 

This is LKE8. 

Here $\fbd$ may have initial twists with $s$ before traversing $i^-\dbar$. On the other hand, $\fbr$ must necessarily continue:
\[
\fbr = s^- i^- \dbar b^*... 
\]
and this is irregardless whether $\fbd$ has any initial twists.

Also, $\fbs$ must begin with $b^+$ and to the left of $\fbp$, otherwise $\fbp$ would trace immediately after crossing $d$. So we have
\[
\fbs = b^+ d o^*...
\]

Now consider $\fbp$, the incoming $\fbr$ forces:
\[
\fbp = b^+ d...
\]

We now have LKE9.

Here $\fbp$ can either continue with $o^*$, or $i^+$. Suppose it were the former. Then it forces $\fbs$ to continue with $o^+$ and to the left of $\fbp$. Otherwise, $\fbp$ would eventually trace. This gives us LKE10.

Here $\fbs$, $\fbr$, and $\fbp$ are stuck in a three-way counterclockwise path mill between $b$ and $o$.

So suppose instead $\fbp$ continued with $i^+$:
\[
\fbp = b^+ d \hat{i}^+ s^+ d r^*...
\]

This is LKE11.

Here suppose $\fbp$ continued with $r^\circ$ or $r^-$, forcing $\fbb$ to continue with $r^-$. This gives LKE14.

Here boyth $\fbi$ and $\fbo$ must continue with $p^-$ and to the right of $\fbb$. Otherwise, $\fbb$ would continue with $p^+$ and then trace. This gives LKE13. Here $\fbp$, $\fbi$, $\fbo$, and $\fbb$ are stuck in a four-way clockwise path mill between $r$ and $p$. This is a contradiction.

So suppose instead that $\fbp$ continued with $r^+$ and to the left of $\fbb$.

This is LKE14.

Here $\fbo$ cannot continue with $p^*$, as that would then force either $\fbo$ or $\fbi$ to eventually trace, whichever possible continuation $\fbi$ and $\fbo$ were to take. So $\fbo$ must continue with $s^-$, but then it will trace, as shown in LKE15.

This is a contradiction.

\subsection{YVQ}\label{YVQ}
Back to the end of Case B,C,D \ref{CaseBCD Track 13}.
Suppose $\fbd$ began with $i$, and $\fbdb$ began with $b$. We remind the reader that we are assuming $\fbd$ must go over an odd number of letters before crossing. This assumption necessitates that eventually $\fbd$ must traverse $i^-$ right before crossing $\dbar$.

We have the situation depicted in YVQ1.

But now at either $\fbo$ or $\fbi$ will immediately trace after crossing $d$

This is a contradiction.

\subsection{SOB}\label{SOB}
Back to the end of Case B,C,D \ref{CaseBCD Track 13}.
Suppose $\fbd$ began with $i$, and $\fbdb$ began with $r$. We remind the reader that we are assuming $\fbd$ must go over an odd number of letters before crossing. This assumption necessitates that eventually $\fbd$ must traverse $i^-$ right before crossing $\dbar$.

We have the situation depicted in SOB1.

Here $\fbd$ must go below $\fbs$ and $\fbp$, or it will be in Situation $\beta$. So we have SOB2.

This then indicates that $\fbd$ cannot have initial twists, as that would cause either $\fbp$ or $\fbs$ to trace. So we have SOB3.

Here the dotted lines are in Situation $\alpha$.

This is a contradiction.

\subsection{BAM}\label{BAM}
Back to the end of Case B,C,D \ref{CaseBCD Track 13}.
Suppose $\fbd$ began with $s$, and $\fbdb$ began with $r$. We remind the reader that we are assuming $\fbd$ must go over an odd number of letters before crossing. This assumption necessitates that eventually $\fbd$ must traverse $s^+$ right before crossing $\dbar$.

Then we must have
\[
\fbr = \dbar b^*...
\]

We have the situation depicted in BAM1. Here it is clear that $\fbp$ will trace.

This is a contradiction.


\subsection{PIZ}\label{PIZ}
Back to the end of Case B,C,D \ref{CaseBCD Track 13}.
Suppose $\fbd$ began with $s$, and $\fbdb$ began with $d$. We remind the reader that we are assuming $\fbd$ must go over an odd number of letters before crossing. This assumption necessitates that eventually $\fbd$ must traverse $s^+$ right before crossing $\dbar$.


We have the situation depicted in PIZ1.

Here the dotted lines must join. Otherwise, either $\fbd$ or $\fbdb$ would immediately be in Situation $\beta$.

This gives PIZ2.

Here $\fbp$ can continue with $o^\circ$, $o^-$, $o^+$, or $i^+$. Whereas $i^\cric$ and $i^-$ would cause $\fbb$ to begin with $i^-$ and then trace immediately after crossing $\dbar$.

\subsubsection{Nurl}
Suppose $\fbp$ continued with $o^\circ$:
\[
\fbp = b^+ d o^\circ
\]

Then $\fbs$ is forced to continue with $o^+$:
\[
\fbs = b^+ d \hat{o}^+ b^*...
\]

This gives Nurl1.

Here if $\fbr$ continued with $b^\circ$ or $b^-$ then $\fbs$ will trace immediately after crossing $d$. But then if $\fbr$ continued with $b^+$ and to the left of $\fbs$, then it would be absorbed:
\[
\fbr = \dbar \hat{b}^+ d o^- \dbar b^- \Abso{\fbp}{\fbs}
\]

This is a contradiction.

\subsubsection{Oast}
Back to the end of PIZ \ref{PIZ}. Suppose $\fbp$ continued with $o^-$:
\[
\fbp = b^+ d o^-...
\]
Then it continues:
\[
\fbp = b^+ d o^- \dbar b^-...
\]

This is Oast1.

Here $\fbp$ and $\fbi$ must continue with $r^-$ and to the right of $\fbo$. Otherwise, $\fbo$ would begin with $r^+$ and immediately trace after crossing $d$.

Then $\fbp$ continues with $s^-$ after crossing, forcing $\fbi$ to do the same. But now $\fbi$ must trace, irregardless of whether $\fbd$ has any initial twists, as shown in Oast2.

This is a contradiction.

\subsubsection{Pilf}
Back to the end of PIZ \ref{PIZ}. Suppose $\fbp$ continued with $o^+$:
\[
\fbp = b^+ d o^+...
\]

This is Pilf1.

Here $\fbp$ must continue with $\dbar$ and not $s^*$, otherwise $\fbs$ would trace. Thus we have Pilf2.

Here similarly $\fbs$ must continue with $\dbar$ to give Pilf3.

Here $\fbr$ must continue with $b^+$ and to the left of $\fbp$ and $\fbs$. Otherwise, $\fbs$ would trace immediately after crossing.

This is Pilf4.

Here we claim that $\fbp$, $\fbs$, and $\fbr$ are stuck in a three-way counterclockwise path mill between $b$ and $o$. Indeed, suppose instead that here $\fbp$ continued with $b^-$ and threby escape to the outer perimeter. Then we note that $\fbb$ must begin with $i^+$ or $i^\circ$, and so $\fbp$ continues:
\[
\fbp  b^- d o^+ s^+ \dbar r^*...
\]

to give us Pilf5. This is regardless whether $\fbd$ has any initial twists or not. Now here if $|fbo$ were to continue with $p^*$, then $\fbi$ will be forced to continue with $p^-$ and to the right of the incoming $\fbo$, otherwise $|fbo$ will eventually trace. This is Pilf6. 

But here it is clear that $\fbp$, $\fbi$, and $\fbo$ are stuck in a clockwise three-way path mill between $r$ and $p$.

So instead suppose $|fbo$ continued with $s^-=$:
\[
\fbo = ...s^- i^- \dbar b^*...
\]

This is shown in Pilf7.

Here it is clear that $\fbs$, $\fbo$ ,and $\fbr$ are stuck in a three-way counterclockwise path mill between $b$ and $o$. Hence, it must not be possible for $\fbp$ to escape to the outer perimeter.

Hence the situation depicted in Pilf4 is indeed a path mill, as it is clear that $\fbs$ could not escape to the outer perimeter or it will trace, and similarly for $r$; and we just showed this as well for $p$.

This is a contradiction.

\subsubsection{Risp}
Back to the end of PIZ \ref{PIZ}. Suppose $\fbp$ continued with $i^+$:
\[
\fbp = b^+ d ^+...
\]

This is Risp1.

Here $\fbi$ must continue with
\[
\fbi = r^- d p^*...
\]
and the incoming $\fbp$ forces
\[
\fbo = r^- d...
\]

as shown in Risp2.

Here $\fbo$ can continue with $p^*$ or $s^-$. Suppose it were the former. Then this forces $\fbi$ to continue with $p^-$ and to the right of the incoming $\fbo$. Otherwise $\fbo$ would be forced to continue with $p^+$ and eventually trace. 

This is Risp3.

Here $\fbp$, $\fbi$, and $\fbo$ are struck in a three-way clockwise path mill between $r$ and $p$. 

So suppose instead $\fbo$ continued with $s^-$. But now it will be absorbed:
\[
\fbo = r^- d \hat{s}^- i^- \dbar b^- \Abso{\fbp}{\fbs}
\]
as shown in Risp4.

This is a contradiction.

\subsection{RUT}\label{RUT}
Back to the end of Case B,C,D \ref{CaseBCD Track 13}.
Suppose $\fbd$ began with $s$, and $\fbdb$ began with $b$. We remind the reader that we are assuming $\fbd$ must go over an odd number of letters before crossing. This assumption necessitates that eventually $\fbd$ must traverse $s^+$ right before crossing $\dbar$.


We have the situation depicted in RUT1.

Here $\fbdb$ must begin with $b^+$ and to the left of $\fbr$. Otherwise $\fbr$ would immediately trace.

This is RUT2.

Here $\fbd$ must go above $\fbo$ and $\fbi$ or it will be in Situation $\beta$. 

This is RUT3.

Here $\fbd$ must continue with $r^+$ and to the left of $\fbp$ and $\fbs$. Otherwise, $\fbp$ or $\fbs$ will trace immediately after crossing $d$. 

This gives us RUT4.

Here we see that the dotted lines are in Situation $\alpha$.

This is a contradiction.

\subsection{ART}\label{ART}
Back to the end of Case B,C,D \ref{CaseBCD Track 13}.
Suppose $\fbd$ began with $p$, and $\fbdb$ began with $r$. We remind the reader that we are assuming $\fbd$ must go over an odd number of letters before crossing. This assumption necessitates that eventually $\fbd$ must traverse $p^+$ right before crossing $\dbar$.


We have the situation depicted in ART1.

Here $\fbp$ will trace immediately after crossing $d$.

This is a contradiction.

\subsection{WOT}\label{WOT}
Back to the end of Case B,C,D \ref{CaseBCD Track 13}.
Suppose $\fbd$ began with $p$, and $\fbdb$ began with $b$. We remind the reader that we are assuming $\fbd$ must go over an odd number of letters before crossing. This assumption necessitates that eventually $\fbd$ must traverse $p^+$ right before crossing $\dbar$.


We have the situation depicted in WOT1.

Here $\fbs$ must begin with $r^+$ otherwise $\fbp$ would trace immediately after crossing $d$. Then $\fbs$ must continue with $b^+$ otherwise $\fbr$ would immediately trace. So we have
\[
\fbs = r^+ b^+ d...
\]

This is WOT2.

Here $\fbd$ must go above $\fbo$ and $\fbi$. Otherwise, it would immediately be in Situation $\beta$. It then continues:L
\[
\fbdb = ...\dbar r^+ b^+ d...
\]

This is WOT3.

Now consider $\fbdb$. It must begin with $b^+$ and to the left of $\fbr$. Otherwise, $\fbr$ would immediately trace. So we have
\[
\fbdb = b^+ d...
\]

This is WOT4. 

Here the dotted lines are in Situation $\alpha$.

This is a contradiction.


\subsection{ORE}\label{ORE}
Back to the end of Case B,C,D \ref{CaseBCD Track 13}.
Suppose $\fbd$ began with $p$, and $\fbdb$ began with $d$. We remind the reader that we are assuming $\fbd$ must go over an odd number of letters before crossing. This assumption necessitates that eventually $\fbd$ must traverse $p^+$ right before crossing $\dbar$.


We have the situation depicted in ORE1.

Here both $\fbp$ and $\fbs$ must begin with $b^+$ and to the left of the incoming $\fbr$. Otherwise, $\fbr$ would immediately trace.

This is ORE2.

Here the dotted lines must join. Otherwise, either $\fbd$ or $\fbdb$ will immediately enter into Situation $\beta$.

So we have ORE3.

Here $\fbp$ can continue with $o^*$ or $i^*$. Suppose it were the former case. Then $\fbs$ also continues with $o^*$. This means that $\fbi$ must begin with $r^-$ and to the right of $\fbo$. Otherwise $\fbo$ would trace:
\[
\fbo = r^+ b^+ d o^\times
\]

So we must have
\[
\fbi = r^- d p^- o^*...
\]
as shown in ORE4.

Here $\fbb$ is forced to begin with $o^*$, as shown in ORE5.

Now if $\fbb$ continued with $o^\circ$ or $o^+$ then $\fbp$ or $\fbs$ would continue with $o^+$ and one of them will immediately trace. But if $\fbb$ instead continued with $o^-$ it would immediately trace after crossing $\dbar$.

So suppose instead $\fbp$ continued with $i^*$, then it must in fact be $i^+$, otherwise $\fbb$ would begin with $i^-$ and then trace immediately after crossing $\dbar$. Then it must continue with $s^*$ to not trace:
\[
\fbp = b^+ d i^+ s^*...
\]

This forces $\fbs$ to continue with $o^*$ to not trace:
\[
\fbs = b^+ d o^*...
\]

Now consider $\fbi$. If it began with $r^\circ$ or $r^+$ then it would force $\fbo$ to begin with $r^+$ and then eventually trace due to the position of $\fbs$:
\[
\fbo = r^+ b^+ d o^\times
\]

Thus we must have
\[
\fbi = r^-...
\]

but then it will trace:
\[
\fbi = r^- d p^- i^\times
\]

This is a contradiction.

This completes the case analysis for when we assume $\fbd$ must go over an odd number of letters before crossing.

We now assume triangle are swapped, and that $\fbd$ must go over an even number of letters before crossing, which necessitates that $\fbdb$ must go over an odd number of letters before crossing.

The possible places to start for $\fbd$ are $p$, $s$, $i$, and $o$; while $\fbdb$ can only start at $b$ or $r$.


\subsection{BTV}\label{BTV}
Back to the end of Case B,C,D \ref{CaseBCD Track 13}.
Suppose $\fbd$ began with $p$, and $\fbdb$ began with $r$. We assume $\fbd$ must go over an even number of letters before crossing, which necessitates that $\fbdb$ must go over an odd number of letters before crossing.

We have the situation depicted in BTV1.

Here due to the assumption that $\fbd$ go over an even number of letters before crossing, it must begin with $p^-$. But now it is clear that $\fbp$ will trace.

This is a contradiction.

\subsection{AFN}\label{AFN}
Back to the end of Case B,C,D \ref{CaseBCD Track 13}.
Suppose $\fbd$ began with $p$, and $\fbdb$ began with $b$. We assume $\fbd$ must go over an even number of letters before crossing, which necessitates that $\fbdb$ must go over an odd number of letters before crossing.

We have the situation depicted in AFN1.

Here due to the assumption that $\fbd$ go over an even number of letters before crossing, it must begin with $p^-$. But now it is clear that either $\fbo$ or $\fbi$ will trace.

This is a contradiction.

\subsection{CSS}\label{CSS}
Back to the end of Case B,C,D \ref{CaseBCD Track 13}.
Suppose $\fbd$ began with $s$, and $\fbdb$ began with $r$. We assume $\fbd$ must go over an even number of letters before crossing, which necessitates that $\fbdb$ must go over an odd number of letters before crossing.

We have the situation depicted in CSS1.

Here it is clear that wither $\fbp$ or $\fbs$ will immediately trace after crossing $d$.

This is a contradiction.

\subsection{UUU}\label{UUU}
Back to the end of Case B,C,D \ref{CaseBCD Track 13}.
Suppose $\fbd$ began with $s$, and $\fbdb$ began with $b$. We assume $\fbd$ must go over an even number of letters before crossing, which necessitates that $\fbdb$ must go over an odd number of letters before crossing.

This is UUU1.

Here $\fbs$ must begin with $r^+$ and to the left of $\fbp$. Otherwise, $\fbp$ would force $\fbo$ and $\fbi$ to cross and then continue with $s^-$, and one of them will trace, as shown in UUU2.

Also, $\fbdb$ must begin with $b^+$ and to the left of $\fbr$ or $\fbr$ will immediately trace. This gives us UUU3.

Here $\fbdb$ can continue with $o^*$ or $i^*$. Suppose it were the former, then it must continue with $o^+ s^*...$, as to not violate Rule F1. But then $\fbs$ will trace, as shown in UUU4.

So suppose instead $\fbdb$ continued with $i^+$. This is UUU5. 

Here if the dotted lines do not join, then it must do so with $\fbdb$ to the right of $\fbd$, as to not put $\fbdb$ in Situation $\beta$, as in UUU6. But now we see $\fbdb$ still enters Situation $\beta$ on the left hand side.

So the dotted lines must join, as shown in UUU7.

Now $\fbo$ can either continue with $p^*$, or $s^-$. Suppose it were the former,then $\fbi$ is forced to continue with $p^-$ otherwise $\fbo$ would trace after crossing $d$. This gives UUU8. Here $\fbr$, $\fbi$, and $\fbo$ are stuck in a three-way clockwise path mill between $b$ and $p$.

So suppose $\fbo$ continued with $s^-$. Then it continues:
\[
\fbo = d s^- i^- \dbar b^- r^*...
\]
forcing
\[
\fbp = r^+ b^+ d...
\]
as shown in UUU9.

Here $\fbp$ can continue with $o$ or $i^+$. Suppose it were the former. Then this forces $\fbs$ to continue with $o^+$ and to the left of the incoming $\fbp$. Otherwise, $\fbp$ would continue with $o^-$ and then trace:
\[
\fbp = ...o^- \dbar b^- r^- d p^\times
\]

So we have
\[
\fbs = o^+ \dbar b^- r^*...
\]

This is UUU10.

Here if $\fbs$ were to continue with $r^\circ$ or $r^+$, then $\fbo$ would continue with $r^+$ and to the left of $\fbs$. In which case, it would trace:
\[
\fbo = ...r^+ b^+ d o^\times
\]
However, if instead $\fbs$ continued with $r^-$ and to the right of $\fbo$, then it would trace:
\[
\fbs = r^- b^+ d i^+ s^\times
\]

So suppose $\fbp$ continued with $i^+$:
\[
\fbp = r^+ b^+ d \hat{i}^+ s^+ \Abso{\fbi}{\fbo}
\]
as shown in UUU11.

This is a contradiction.

\subsection{GTK}\label{GTK}
Back to the end of Case B,C,D \ref{CaseBCD Track 13}.
Suppose $\fbd$ began with $i$, and $\fbdb$ began with $b$. We assume $\fbd$ must go over an even number of letters before crossing, which necessitates that $\fbdb$ must go over an odd number of letters before crossing.

This is GTK1..

But here $\fbb$ must continue with $r^\circ$ or $r^-$ otherwise it would immediately trace. However, this means 
\[
\fbp = r^- d...
\]
and
\[
\fbs = r^- d...
\]
and one of these will immediately trace.

This is a contradiction.


\subsection{CGI}\label{CGI}
Back to the end of Case B,C,D \ref{CaseBCD Track 13}.
Suppose $\fbd$ began with $i$, and $\fbdb$ began with $r$. We assume $\fbd$ must go over an even number of letters before crossing, which necessitates that $\fbdb$ must go over an odd number of letters before crossing.

This is CGI1.

Here $\fbdb$ must begin with $r^-$ and to the right of $\fbb$, otherwise $\fbb$ would immediately trace. This gives us CGI2.

Here $\fbd$ can continue with either $s^+ \dbar...$, or $p^+ \dbar...$ One can check that there are no other possibilities without violating the assumption that it must go over an even number of letters before crossing.

But in both cases, this forces 
\[
\fbo = b^+ d o^\times
\]

as if it were to begin with $b^\circ$ or $b^-$ then $\fbi$ would begin with $b^-$ and then eventually trace.

This is shown in CGI3, and CGI4.

\subsection{GLD}\label{GLD}
Back to the end of Case B,C,D \ref{CaseBCD Track 13}.
Suppose $\fbd$ began with $o$, and $\fbdb$ began with $b$. We assume $\fbd$ must go over an even number of letters before crossing, which necessitates that $\fbdb$ must go over an odd number of letters before crossing.

This is GLD1.

But here it is immediately clear that $\fbo$ will trace.


\subsection{UGN}\label{UGN}
Back to the end of Case B,C,D \ref{CaseBCD Track 13}.
Suppose $\fbd$ began with $o$, and $\fbdb$ began with $r$. We assume $\fbd$ must go over an even number of letters before crossing, which necessitates that $\fbdb$ must go over an odd number of letters before crossing.

This is UGN1.

Here we must have
\[
\fbp  = d o^+...
\]
and
\[
\fbs = d o^+...
\]

But now one of them will immediately trace.


This concludes the case analyses for Track 13.
